\begin{examplebox}
	\textbf{\textcolor{CExample}{例 1（基础）}}

	\textbf{\textcolor{CExample}{题目：}}
	已知椭圆
	\[
		E:\frac{x^2}{4}+y^2=1,
	\]
	直线
	\[
		l:3x-4y+5=0.
	\]
	求椭圆 \(E\) 上的点到直线 \(l\) 的距离的最小值与最大值。

	\textbf{\textcolor{CExample}{解题步骤：}}

	\begin{description}[style=nextline, leftmargin=2.8em, labelsep=0.5em, itemsep=0.45em]
		\item[\textbf{第一步（计算基本量）：}]
		      由椭圆方程得
		      \[
			      a=2,\qquad b=1.
		      \]
		      由直线方程得
		      \[
			      A=3,\qquad B=-4,\qquad C=5.
		      \]
		      所以
		      \[
			      R=\sqrt{a^2A^2+b^2B^2}
			      =
			      \sqrt{4\times 9+1\times 16}
			      =
			      \sqrt{52}
			      =
			      2\sqrt{13},
		      \]
		      \[
			      \sqrt{A^2+B^2}=5.
		      \]

		      \par\textbf{理由：}
		      代入公式前必须计算出 \(R\) 和分母，它们是后续比较和代入的基础。

		\item[\textbf{第二步（比较 \(|C|\) 与 \(R\)：}]
		      \[
			      |C|=5,
			      \qquad
			      R=2\sqrt{13}\approx 7.21.
		      \]
		      因为
		      \[
			      |C|<R,
		      \]
		      所以直线与椭圆相交。

		      \par\textbf{理由：}
		      相交时最小距离为 \(0\)；若相离，则最小距离才需要使用 \(|C|-R\)。

		\item[\textbf{第三步（套用最值公式）：}]
		      由 \(|C|\le R\) 得
		      \[
			      d_{\min}=0,
		      \]
		      \[
			      d_{\max}
			      =
			      \frac{|C|+R}{\sqrt{A^2+B^2}}
			      =
			      \frac{5+2\sqrt{13}}{5}.
		      \]

		      \par\textbf{理由：}
		      最小值直接为 \(0\)，最大值代入公式即得。
	\end{description}

	\textbf{\textcolor{CExample}{关键结论：}}
	\[
		d_{\min}=0,
		\qquad
		d_{\max}=\frac{5+2\sqrt{13}}{5}.
	\]

	\medskip

	\textbf{\textcolor{CExample}{例 2（稍变形）}}

	\textbf{\textcolor{CExample}{题目：}}
	已知椭圆
	\[
		E:\frac{x^2}{9}+\frac{y^2}{4}=1,
	\]
	直线
	\[
		l:4x+3y+20=0.
	\]
	求椭圆 \(E\) 上的点到直线 \(l\) 的距离的最小值与最大值。

	\textbf{\textcolor{CExample}{解题步骤：}}

	\begin{description}[style=nextline, leftmargin=2.8em, labelsep=0.5em, itemsep=0.45em]
		\item[\textbf{第一步（计算基本量）：}]
		      由题意得
		      \[
			      a=3,\qquad b=2,
		      \]
		      \[
			      A=4,\qquad B=3,\qquad C=20.
		      \]
		      所以
		      \[
			      R=\sqrt{a^2A^2+b^2B^2}
			      =
			      \sqrt{9\times 16+4\times 9}
			      =
			      \sqrt{144+36}
			      =
			      \sqrt{180}
			      =
			      6\sqrt{5},
		      \]
		      \[
			      \sqrt{A^2+B^2}=5.
		      \]

		      \par\textbf{理由：}
		      先计算 \(R\) 和分母，为后续步骤做准备。

		\item[\textbf{第二步（比较 \(|C|\) 与 \(R\)：}]
		      \[
			      |C|=20,
			      \qquad
			      R=6\sqrt5\approx 13.42.
		      \]
		      因为
		      \[
			      |C|>R,
		      \]
		      所以直线与椭圆相离。

		      \par\textbf{理由：}
		      相离时最小距离为正，需用 \(|C|-R\) 计算。

		\item[\textbf{第三步（套用最值公式）：}]
		      \[
			      d_{\min}
			      =
			      \frac{|C|-R}{\sqrt{A^2+B^2}}
			      =
			      \frac{20-6\sqrt5}{5},
		      \]
		      \[
			      d_{\max}
			      =
			      \frac{|C|+R}{\sqrt{A^2+B^2}}
			      =
			      \frac{20+6\sqrt5}{5}.
		      \]

		      \par\textbf{理由：}
		      由 \(|C|>R\)，最小值和最大值分别对应 \(|C|-R\) 和 \(|C|+R\) 除以分母。
	\end{description}

	\textbf{\textcolor{CExample}{关键结论：}}
	\[
		d_{\min}=\frac{20-6\sqrt5}{5},
		\qquad
		d_{\max}=\frac{20+6\sqrt5}{5}.
	\]

	\medskip

	\textbf{\textcolor{CExample}{例 3（综合应用）}}

	\textbf{\textcolor{CExample}{题目：}}
	已知椭圆
	\[
		E:\frac{x^2}{4}+\frac{y^2}{b^2}=1\quad (0<b<2)
	\]
	与直线
	\[
		l:3x-4y+5=0.
	\]
	椭圆 \(E\) 上的点到 \(l\) 的距离的最小值为 \(0\)，最大值为
	\[
		\frac{13}{5}.
	\]
	求实数 \(b\) 的值。

	\textbf{\textcolor{CExample}{解题步骤：}}

	\begin{description}[style=nextline, leftmargin=2.8em, labelsep=0.5em, itemsep=0.45em]
		\item[\textbf{第一步（写出 \(R\) 表达式）：}]
		      由题意得
		      \[
			      a=2,
			      \qquad
			      A=3,\qquad B=-4,\qquad C=5.
		      \]
		      所以
		      \[
			      R=\sqrt{a^2A^2+b^2B^2}
			      =
			      \sqrt{4\times 9+b^2\times 16}
			      =
			      \sqrt{36+16b^2},
		      \]
		      \[
			      \sqrt{A^2+B^2}=5.
		      \]

		      \par\textbf{理由：}
		      本题 \(b\) 未知，\(R\) 含 \(b\)，后续需利用最大距离条件列方程。

		\item[\textbf{第二步（利用 \(d_{\min}=0\) 检查位置关系）：}]
		      最小值为 \(0\)，说明直线与椭圆有公共点，即
		      \[
			      |C|\le R.
		      \]
		      这里
		      \[
			      |C|=5,
			      \qquad
			      R=\sqrt{36+16b^2}.
		      \]
		      由于 \(b>0\)，所以
		      \[
			      R=\sqrt{36+16b^2}>6>5.
		      \]
		      因此 \(|C|\le R\) 自动满足，\(d_{\min}=0\) 不再额外限制 \(b\)。

		      \par\textbf{理由：}
		      本题真正决定 \(b\) 的条件是最大距离，而不是最小距离。

		\item[\textbf{第三步（利用 \(d_{\max}=\frac{13}{5}\) 列方程）：}]
		      由最大值公式
		      \[
			      d_{\max}
			      =
			      \frac{|C|+R}{\sqrt{A^2+B^2}},
		      \]
		      得
		      \[
			      \frac{5+\sqrt{36+16b^2}}{5}
			      =
			      \frac{13}{5}.
		      \]
		      于是
		      \[
			      5+\sqrt{36+16b^2}=13,
		      \]
		      \[
			      \sqrt{36+16b^2}=8.
		      \]
		      两边平方得
		      \[
			      36+16b^2=64,
		      \]
		      所以
		      \[
			      16b^2=28,
			      \qquad
			      b^2=\frac{7}{4}.
		      \]
		      由于 \(b>0\)，所以
		      \[
			      b=\frac{\sqrt7}{2}.
		      \]

		      \par\textbf{理由：}
		      这样解出的 \(b=\frac{\sqrt7}{2}\) 满足 \(0<b<2\)，因此仍是长轴在 \(x\) 轴上的椭圆，与前面的使用条件一致。
	\end{description}

	\textbf{\textcolor{CExample}{关键结论：}}
	\[
		b=\frac{\sqrt7}{2}.
	\]

\end{examplebox}